\documentclass[a4paper, 10pt, ]{article}

\usepackage[slovak]{babel}

% ------------------------------

\usepackage[utf8]{inputenc}
\usepackage[T1]{fontenc}

\usepackage[left=4cm,
            right=4cm,
            top=2.1cm,
            bottom=2.6cm,
            footskip=7.5mm,
            twoside,
            marginparwidth=3.0cm,
            %showframe,
            ]{geometry}

\usepackage{graphicx}
\usepackage[dvipsnames]{xcolor}
% https://en.wikibooks.org/wiki/LaTeX/Colors

% ------------------------------

\usepackage{lmodern}

\usepackage[tt={oldstyle=false,proportional=true,monowidth}]{cfr-lm}
% https://mirror.szerverem.hu/ctan/fonts/cfr-lm/doc/cfr-lm.pdf

% ------------------------------

\usepackage{amsmath}
\usepackage{amssymb}
\usepackage{amsthm}

\usepackage{booktabs}
\usepackage{multirow}
\usepackage{array}
\usepackage{dcolumn}

\usepackage{natbib}

% ------------------------------

\hyphenpenalty=6000
\tolerance=1000

\def\naT{\mathsf{T}}

% ------------------------------

\makeatletter

    \def\@seccntformat#1{\protect\makebox[0pt][r]{\csname the#1\endcsname\hspace{4mm}}}

    \def\cleardoublepage{\clearpage\if@twoside \ifodd\c@page\else
    \hbox{}
    \vspace*{\fill}
    \begin{center}
    \phantom{}
    \end{center}
    \vspace{\fill}
    \thispagestyle{empty}
    \newpage
    \if@twocolumn\hbox{}\newpage\fi\fi\fi}

    \newcommand\figcaption{\def\@captype{figure}\caption}
    \newcommand\tabcaption{\def\@captype{table}\caption}

\makeatother

% ------------------------------

\usepackage{fancyhdr}
\fancypagestyle{plain}{%
\fancyhf{} % clear all header and footer fields
% \fancyfoot[C]{\sffamily {\bfseries \thepage}\ | {\scriptsize\oznacenieCasti}}
\fancyfoot[C]{\sffamily {\bfseries \thepage}{\color{Gray}\scriptsize$\,$z$\,$\pageref{LastPage}}\ | \includegraphics[height=5pt]{../../COMMONFILES/KUT_logo_v0.1.pdf}{\scriptsize\KUTporadoveCislo}}
\renewcommand{\headrulewidth}{0pt}
\renewcommand{\footrulewidth}{0pt}}
\pagestyle{plain}

% ------------------------------

\usepackage{titlesec}
\titleformat{\paragraph}[hang]{\sffamily  \bfseries}{}{0pt}{}
\titlespacing*{\paragraph}{0mm}{3mm}{1mm}
\titlespacing*{\subparagraph}{0mm}{3mm}{1mm}

\titleformat*{\section}{\sffamily\Large\bfseries}
\titleformat*{\subsection}{\sffamily\large\bfseries}
\titleformat*{\subsubsection}{\sffamily\normalsize\bfseries}


% ------------------------------

\PassOptionsToPackage{hyphens}{url}
\usepackage[pdfauthor={},
            pdftitle={},
            pdfsubject={},
            pdfkeywords={},
            % hidelinks,
            colorlinks=false,
            breaklinks,
            ]{hyperref}


% ------------------------------

\graphicspath{%
{../fig_standalone/}%
{../../PY/fig/}%
{../../ML/fig/}%
{./fig/}%
}

% ------------------------------

\usepackage{enumitem}

\usepackage{lettrine}

% ------------------------------

\usepackage{lastpage}

\usepackage{microtype}

% ------------------------------

\usepackage{algorithm}
\usepackage[noend]{algpseudocode}
\makeatletter
\renewcommand{\ALG@name}{Algoritmus}
\makeatother
\usepackage{amsmath}
\usepackage{bbold}
\usepackage{calc}
\usepackage{dsfont}
\usepackage{mathtools}
\usepackage{tabto}


\newcommand{\mr}[1]{\mathrm{#1}}
\newcommand{\bs}[1]{\boldsymbol{#1}}
\newcommand{\bm}[1]{\mathbf{#1}}

\newcommand{\diff}[2]{\frac{\Delta #1}{\Delta #2}}
\newcommand{\der}[2]{\frac{d #1}{d #2}}
\newcommand{\parder}[2]{\frac{\partial #1}{\partial #2}}

\newcommand{\argmax}[0]{\mr{argmax}}
\newcommand{\diag}[0]{\mr{diag}}
\newcommand{\rank}[0]{\mr{rank}}
\newcommand{\trace}[0]{\mr{tr}}

\renewcommand{\Re}{\mr{Re}}
\renewcommand{\Im}{\mr{Im}}


\theoremstyle{definition}
\newtheorem{definition}{Definícia}[section]
\newtheorem{theorem}{Veta}[section]
\newtheorem{lemma}[theorem]{Lemma}
\newtheorem{example}{Príklad}[section]
\renewcommand*{\proofname}{Dôkaz}

% ------------------------------

\usepackage{siunitx}
\usepackage{xcolor}

% Define a custom color
\definecolor{backcolour}{rgb}{0.95,0.95,0.92}
\definecolor{codegreen}{rgb}{0,0.6,0}

% Define a custom style
\lstdefinestyle{myStyle}{
    belowcaptionskip=1\baselineskip,
    breaklines=true,
    showspaces=false,                % show spaces everywhere adding particular underscores; it overrides 'showstringspaces'
    showstringspaces=false,          % underline spaces within strings only
    showtabs=false, 
    frame=single,
    numbers=left,
    basicstyle=\footnotesize\ttfamily,
    keywordstyle=\bfseries\color{blue!70!black},
    commentstyle=\color{green!40!black},
    identifierstyle=\color{gray!50!black},
    backgroundcolor=\color{white},
}

% Use \lstset to make myStyle the global default
\lstset{style=myStyle}

% -----------------------------------------------------------------------------

\def\oznacenieCelku{Kolekcia učebných textov}

% -----------------------------------------------------------------------------


\def\KUTporadoveCislo{devAS1}

\def\oznacenieVerzie{v1.0}
% \def\oznacenieVerzie{\phantom{v1.0}}

\def\mesiacRok{september 2025}

\def\authorslabel{RJ}


% Figures: "obr. 1"
\newcommand{\figref}[1]{\hyperref[#1]{obr.~\ref*{#1}}}

% Tables: "tab. 1"
\newcommand{\tabref}[1]{\hyperref[#1]{tab.~\ref*{#1}}}

% Equations: "rov. (1)"
\newcommand{\eqrefsk}[1]{\hyperref[#1]{rov.~(\ref*{#1})}}



% -----------------------------------------------------------------------------

\begin{document}

% -----------------------------------------------------------------------------
% Uvodny nadpis

\noindent
\parbox[t][18mm][c]{0.3\textwidth}{%
    \raisebox{-0.9\height}{%
        \phantom{.}\includegraphics[height=18mm]{./COMMONFILES/URKFEIlogo.pdf}%
    }%
}%
\parbox[t][18mm][c]{0.7\textwidth}{%
    \raggedleft

    \sffamily
    \fontsize{16pt}{18pt}
    \fontseries{sbc}
    \selectfont

    \noindent
    \textcolor[rgb]{0.75, 0.75, 0.75}{\textls[25]{\oznacenieCelku}}
}%

\noindent
\parbox[t][16mm][b]{0.5\textwidth}{%
    \raggedright

    \color{Gray}
    \sffamily

    \fontsize{12pt}{12pt}
    \selectfont
    \mesiacRok

    \fontsize{6pt}{10pt}
    \selectfont
    \href{https://github.com/OkoliePracovnehoBodu/KUT}{github.com/OkoliePracovnehoBodu/KUT}

    \fontsize{8pt}{10pt}
    \selectfont
    \authorslabel




}%
\parbox[t][16mm][b]{0.5\textwidth}{%
    \raggedleft

    \sffamily

    \fontsize{6pt}{6pt}
    \selectfont

    \textcolor[rgb]{0.68, 0.68, 0.68}{\oznacenieVerzie}


    \fontsize{14pt}{14pt}
    \selectfont

    \bfseries

    \includegraphics[height=12pt]{./COMMONFILES/KUT_logo_v0.1.pdf}%
    {%
        \textls[-50]{\KUTporadoveCislo}
    }%
}%

% -----------------------------------------------------------------------------




\vspace{6mm}

% ---------------------------------------------
\sffamily
\bfseries
\fontsize{18pt}{21pt}
\selectfont

\begin{flushleft}
    Laboratórne zariadenie AeroShield:\\ softwáre Matlab
\end{flushleft}

\bigskip

% -----------------------------------------------------------------------------
\normalsize
\normalfont
% -----------------------------------------------------------------------------

\lstset{style=mystyle}










\noindent
\lettrine[lines=1, nindent=1pt, loversize=0.0]{C}{ieľom}
textu je zoznámenie sa s obslužným Matlab softwárom pre prácu so zariadením \emph{AS}, ktorý slúži na meranie a ovládanie spomínaného zariadenia.


\section{Hlavný program}
\begin{lstlisting}[caption=Zavolanie funkcie merania., label={code:main}, language=Matlab]
% ----------------------------------
% Run the experiment using the following line of code
% ----------------------------------

[t, y, u, potentiometer, yhat, dyhat] = runArduinoPlot();

% ----------------------------------
% ----------------------------------
\end{lstlisting}

Zavolaním funckie \texttt{runArduinoPlot()} sa začne navrhnuté meranie, ktoré následne automaticky skončí, čí uz po vypršaní času merania alebo inou podmienkov, ktorú si vieme zadefinovať samostatne. Definícia ukončovacích podmienok bude spomenutá nižšie. Ako môžeme vidieť na výpise kódu \ref{code:main}, funkcia vracia dáta, ktoré následne možno použiť.
Ktoré dáta vracia si vieme následne vo funkcii zadefinovať, poprípade si vieme dodať, daľšie pre nás dôležité veličiny.

\section{Dátový priečinok}
\begin{lstlisting}[caption=Vytvorenie dátového priečinku., label={code:repo}, language=Matlab]
% ----------------------------------
% Create the data repository
% ----------------------------------
DDIR = "dataRepo";
if ~exist(DDIR, "dir")
    fprintf("Creating the data directory...");
    mkdir(DDIR);
end
% ----------------------------------
% ----------------------------------
\end{lstlisting}

Kód, ktorý je ukázaný vo výpise \ref{code:repo}, slúži na vytvorenie adresáru (priečinku) pre uloženie meraní. Do tohto priečinka sa bežne ukladajú merania av dvoch formátoch \emph{CSV} a \emph{MAT}. \emph{CSV} súbor obsahuje merané veličiny dostatné po sériovej linke, \emph{MAT} súbor obsahuje celý Matlab workspace - všetky definované premenné v Matlab-e.
\emph{Je odporúčané si manuálne vytvoriť nový adresár, do ktorého si budeme kopírovať (CTRL+C & CTRL+V) merania, s ktorými chceme následne pracovať, aby sme vždy mali netknuté dáta z merania v zálohe.}

\section{Definícia premenných}
\begin{lstlisting}[caption=Definícia všetkých potrebných premenných., label={code:vars}, language=Matlab]
% ----------------------------------
% Define all the parameters
% ----------------------------------

% Define time parameters

T_start = 0;

T_sample = 3;      % [ms] <1, 255>

% Define STOP TIME

T_stop = 60.0;     % [sec]

% Define control parameters
U_MAX = 100.0;
U_MIN = 0.0;
Y_SAFETY = 190.0;

% Define PID param
P = 1.0;
I = 0.30;
D = 0.19;

R_WANTED = 140;

% alpha - beta filter
alpha = 0.8;
beta = 0.2;

timer_t = [];
timer_y = [];
timer_yhat = [];
timer_dyhat = [];
timer_u = [];
timer_potentiometer = [];
% ----------------------------------
% ----------------------------------
\end{lstlisting}

Táto časť kódu (\ref{code:vars}) slúži na definíciu všetkých potrebných premenných pre dané meranie, v tomto prípade sme si zadefinovali premenné pre PID regulátor a $\alpha-\beta$ filter, aby sme ukázali možné riešenia implementácie riadenia a pozorovania. Následne v \tabref{tab:quick_desc} budú potrebné premenné pre priebeh merania popísané, ostatné premenné sú špecificky zvolené pre tento poskytnutý náhľadový kód - nebudú potrebne pre napr. meranie dynamických vlastností systému.

\begin{center}

    \vspace{-10pt}

    \tabcaption{Krátky popisok nevyhnutných premenných}
    \label{tab:quick_desc}

    \lstyle

    \begin{tabular*}{\textwidth}{@{ \extracolsep{\fill}} ll|l}
        \toprule
        Premenná & Jednotka & Popis \\
        \midrule
        \texttt{T\_start} & s & Čas začiatku merania. \\
        \texttt{T\_sample} & ms & Perióda vzorkovania v rozsahu 1 až 255. \\
        \texttt{T\_stop} & s & Čas trvania merania. \\
        \texttt{U\_MAX} & \% & Maximálny akčný zásah. \\
        \texttt{U\_MIN} & \% & Minimálny akčný zásah. \\
        \texttt{Y\_SAFETY} & $^\circ$ & Maximálna hodnota výstupnej veličiny. \\
        \texttt{timer\_t} & s & Pole časov v ktorých prebehla komunikácia. \\
        \texttt{timer\_y} & $^\circ$ & Pole výstupov zo zariadenia AeroShield. \\
        \texttt{timer\_u} & \% & Pole vstupov do zariadenia AeroShield. \\
        \texttt{timer\_potentiometer} & \% & Pole hodnôt z potenciometra. \\
        \bottomrule
    \end{tabular*}
\end{center}

\newpage
\section{Vykreslenie dát v reálnom čase}
\begin{lstlisting}[caption=Definícia časovača na vykreslenie meraných dát v reálnom čase., label={code:plot_timer}, language=Matlab]
% ----------------------------------
% Plot the measured data in real time
% ----------------------------------
function plotData()
    persistent hy hr hu;
    
    % persistent example
    
    try
        if isempty(hy) || isempty(hr) || isempty(hu)
            f = figure(9999); clf(f);
            ax = axes(f);
            hold on;
            hy = plot(ax, nan, nan, '.b');
            hr = plot(ax, nan, nan, '.r');
            hu = plot(ax, nan, nan, '.k');

            % example = plot(ax, nan, nan, '.g');
            
            grid minor;
            title("Real-Time System Response");
            xlabel("t [s]");
            ylabel("$\varphi [^\circ]$", "Interpreter","latex");
            legend(ax, "y","ref", "yhat", 'Location', 'southeast');
            
        end

        set(hy, 'YData', timer_y, 'XData', timer_t);
        set(hr, 'YData', timer_potentiometer, 'XData', timer_t);
        set(hu, 'YData', timer_yhat, 'XData', timer_t);
        
        % set(example, 'YData', example_y, 'XData', example_x);
        
        drawnow limitrate nocallbacks;
    catch err
       fprintf(2, "Plot thread: " + err.message + "\n");
    end
end

tPlot = timer('ExecutionMode','fixedRate', 'Period', 0.5, 'TimerFcn', @(~, ~) plotData());
start(tPlot);
% ----------------------------------
% ----------------------------------
\end{lstlisting}

Obslužná funkcia \texttt{plotData} na vykreslovanie grafov v reálnom čase, ktorá je následne volaná časovačom \texttt{tPlot} v intervale $500$ milisekúnd pokiaľ sa neukončí meranie. V prípade, že by sme chceli dodať daľší signál do grafu, môžeme postupovať podľa už predrobeného príkladu, kedy sme v bloku kódu \ref{code:plot_timer} zakomentovali $3$ riadky. Avšak, keď pridáme nový signál na vykreslovanie, tak je pre funkčnosť a plynulosť dôležité zvýšiť periódu vykreslovania. Nakoľko počet dát na vykreslenie sa zvýši a vykreslovanie v Matlab-e, nie je obzvlásť rýchle.

\section{Záznamové súbory merania}
\begin{lstlisting}[caption=Inicializácia záznamových súborov., label={code:files}, language=Matlab]
% ----------------------------------
% Initialize File Streams
% ----------------------------------

DateString = convertCharsToStrings(datestr(datetime('now'), "yyyy_mm_dd_HH_MM_ss"));

FILENAME = "dataFile";

function fullpath = getfilename(dirpath, filename, datestr, ext)
    if nargin < 3
        error("At least the first 3 parameters need to be provided.");
    elseif nargin == 3
        ext = "csv";
    end

    fullpath = "./" + dirpath + "/" + filename + "_" + datestr + "." + ext;
end

FILEPATH = getfilename(DDIR, FILENAME, DateString);
FILEPATH_MAT = getfilename(DDIR, FILENAME, DateString, 'mat');

if(exist("datafileID", "var"))
    fclose(datafileID);
    clear datafileID;
end

datafileID = fopen(FILEPATH,'w');
fprintf(datafileID, 't, tp, r, y, u, dtp, dt\n');
% ----------------------------------
% ----------------------------------
\end{lstlisting}

Vytvorenie súboru na zapísanie merania a zapísanie názvov stĺpcov, ktoré sú do súboru zapisované. V prípade, že súbor už existuje alebo je aktívny v Matlab prostredí, tak sa zavrie, aby každé meranie bolo vložené do vlastného súboru. Ak chceme zmeniť cestu ukladania súborov merania, tak vo funkcii \texttt{getfilename} musíme zmeniť premennú \texttt{fullpath}. Keď pridáme nový signál do súboru, je nutné pridať aj meno stĺpca, aby nedošlo k nesprávnemu načítaniu, keď budeme dáta potrebovať načítať.

\section{Zapisovanie dát merania}
\begin{lstlisting}[caption=Zapisovanie meraných dát do súboru a konzoly., label={code:files_write}, language=Matlab]
% ----------------------------------
% Write data into files
% ----------------------------------

function updateInfo(datafileID, dt, Ts, x)
    if ((dt) > (Ts*1.05))
        fprintf('%8.3f %8.3f %8.3f %8.3f %8.3f %8.3f %8.3f  --\n', x); % write to console, measurement took too long
    else
        fprintf('%8.3f %8.3f %8.3f %8.3f %8.3f %8.3f %8.3f\n', x); % write to console
    end
    fprintf(datafileID, '%8.3f,%8.3f,%8.3f,%8.3f,%8.3f,%8.3f,%8.3f\n', x); % write to file

    timer_t = [timer_t x(1)];
    timer_y = [timer_y x(4)];
    timer_u = [timer_u x(5)];
    timer_potentiometer = [timer_potentiometer x(3)];
end

doUpdate = @(x) updateInfo(datafileID, x(end), T_sample, x);

% ----------------------------------
% ----------------------------------
\end{lstlisting}

Funkcia \texttt{updateInfo} vo výpise kódu \ref{code:files_write} slúži na zapisovanie dát do súboru \emph{CSV} a do konzoly.

Ak pridáme nový signál do zapisovania dát, musíme pridať \texttt{,\%8.3f} do riadku s komentárom \emph{\% write to file} pred text \texttt{\textbackslash{n}}. Následne sa do súboru bude ukladať nový signál.

\section{Sériová komunikácia}
\begin{lstlisting}[caption=Inicializácia sériovej komunikácie a konfigurácia., label={code:init_serial}, language=Matlab]
% ----------------------------------
% Define serial port parameters, open and configure comms
% ----------------------------------

if(exist("serPort", "var"))
    serPort.flush("input");
    clear serPort;
end

serPort = serialport('COM3', 115200, 'Timeout', 5);

serLine = readline(serPort);

while(~contains(serLine, "config"))
    disp(serLine);
    serLine = readline(serPort);
end

fprintf("Sending now\n");
write(serPort, cast(T_sample, "uint8"), "uint8");

% Read the first line from the serial port (MCU starting)
while(~contains(serLine, "start"))
    disp(serLine);
    serLine = readline(serPort);
end

disp(serLine);
write(serPort, 0.0, 'single'); % Necessary to send this command for stable sampling period

while(contains(serLine, "---"))
    disp(serLine);
    serLine = readline(serPort);
end

% Read and parse the calibration data
serLineList = str2num(serLine); %#ok<ST2NM>

% ----------------------------------
% ----------------------------------
\end{lstlisting}

V tejto sekcii kódu \ref{code:init_serial} popisujeme postup inicializácie sériovej komunikácie s Arduinom a zároveň nastavenie periódy vzorkovania. Je v záujme čitateľa predom porozumieť, fungovanie sériovej komunikácie a poskytnutého kódu pred snahou modifikácie tejto časti kódu.
V skrate, otvárame sériovú komunikácie na porte \texttt{COM3}, v prípade, že dojde k situácii, kedy Matlab vypíše chybné hlásenie o nenajdení alebo nedostupnosti zariadenia na tomto porte, tak sú 3 možnosti, ktoré vieme aplikovať na vyriešenie tohto problému.

\begin{enumerate}
    \item Vypojiť a následne zapojiť zariadenie z počítača pomocou USB káblu.
    \item Pozrieť pripojené fyzické porty k počítaču v termináli pomocou príkazu \texttt{mode} na zariadení s operačným systémom Windows, \texttt{sudo dmesg | grep -i tty} na Linux/MacOS. Ak, nájdeme pripojené zariadenia, tak v prípade Windows-u zmeníme v kóde \texttt{COM3} na iný \texttt{COM#}, pri Linux-e/MacOS je názov portu v tvare \texttt{/dev/ttyS###}, následne aplikujeme rovnaký postup ako pri Windows zariadení.
    \item Nakoniec, môže dojsť ku chybe komunikácie a zariadenie nebude reagovať, tak do Matlab konzoly napíšeme \texttt{save; clear;} a spustíme - uložíme si týmto spôsobom dáta a zaroveň vymažeme celé prostredie Matlab-u. Následne spustíme meranie, ktoré predčasne ukončíme, keď sa inicializuje komunikácia a vstup do systému bude $0$ - predčasné ukočenie vieme dosiahnuť tak, že pri spúštaní podržíme vrtulku v nulovej pozícii (kyvadlo je vo zvyslej polohe, kde má uhoľ natočenia rovný $0$). Keď meranie začne, tak vrtulku manuálne vyvrhneme nad bezpečnostnú hranicu výstupu - nami definovaná hodnota je $190^\circ$. Ak to nepomôže, tak vtedy je potrebné zariadenie odpojiť od napájania (zásuvky) a následne pripojiť, aby sa resetovalo Arduino, podľa možností možeme natvrdo resetovať Arduino.
\end{enumerate}

\section{Počiatočné hodnoty}
\begin{lstlisting}[caption=Zaznamenanie počiatočných hodnôt., label={code:init_values}, language=Matlab]
% ----------------------------------
% Extract the initial values from the received data
% ----------------------------------
plant_time_init = serLineList(1);
plant_potentiometer_init = serLineList(2);
plant_output_init = serLineList(3);
plant_input_init = serLineList(4);

plant_time = serLineList(1) - plant_time_init;
plant_input = serLineList(2);
plant_output = serLineList(3);
plant_potentiometer = R_WANTED + serLineList(4)/100*20;
plant_dt = serLineList(5);

timer_yhat = [timer_yhat, plant_output];
timer_dyhat = [timer_dyhat, 0];

% Display the received data
tmp_printlist = [0, plant_time, plant_potentiometer, plant_output, plant_input, plant_dt, T_sample];
doUpdate(tmp_printlist);
% ----------------------------------
% ----------------------------------
\end{lstlisting}

Kód zobrazený na \ref{code:init_values} sa týka počiatočnej komunikácie so zariadením, počiatočne snímané hodnoty sú cez sériovú linku poslané Arduinom, po čom započne samotné meranie. Počiatočné hodnoty sa používajú v prípade času, kedy odpočítavame od času trvania merania počiatočný čas, aby sme v rámci záznamu začínali časom v $0$, nakoľko čas pred týmto odsekom kódu náleží inicializácii zariadenia a nie samotnému meraniu.
Po prvýkrát voláme funkciu \texttt{doUpdate(x)}, ktorá zaznamená merané dáta do súboru a konzoly. V našom prípade nepoužívame hodnotu potenciometra navrátenú zariadením priamo, nakoľko tento signál používame na jemné menenie žiadanej hodnoty. Nakoľko sa pohybujeme v relatívne nestabilnom pracovnom okolí s našou požiadavkou referenčnej hodnoty, tak sme museli možnosť zmeny referenčnej hodnoty zjemniť natoľko, aby nedošlo k takej zmene, ktorá by navrhnuté riadenie destabilizovalo.

\section{Definícia premenných hlavne slučky}
\begin{lstlisting}[caption=Nastavenie premenných v hlavnej slučke., label={code:loop_vars}, language=Matlab]
% ----------------------------------
% Set the main loop parameters
% ----------------------------------

% Set initial control input value
e_old = 0;
e_int_old = 0;
u = 0;
u_send = u;


% Get the initial time
time_start = datetime('now');
time_tick = time_start;

% ----------------------------------
% ----------------------------------
\end{lstlisting}

Jednoduchá definícia premenných (viď \ref{code:loop_vars}), ktoré sa používaju najmä v rámci slučky merania, sú tu umiestnené čisto z dôvodu, aby sme sa nemuseli hýbať v kóde až na začiatok, keď chceme zmeniť alebo doplniť premenné. Možno sem umiestniť väčšinu premenných, až na tie od ktorých chceme, aby boli zaznamenané v súboroch.

\section{Čítanie sériovej komunikácie}
\begin{lstlisting}[caption=Definícia počúvateľa sériovej komunikácie., label={code:read}, language=Matlab]
% ----------------------------------
% Define serial link listener
% ----------------------------------

function readSerialData(src, ~) 
    data = readline(src);
    src.UserData = data;
end

configureCallback(serPort, "terminator", @readSerialData);
% ----------------------------------
\end{lstlisting}

Výpis \ref{code:read} slúži na zadefinovanie obslužnej funkcie pre čítanie dát zo seriového portu, Matlab ponúka možnosť definície funkcie, ktorá bude volaná, ked zaznamená nové dáta v sériovej zbernici. Funkcionalitu následne používame v slučke na riadenie komunikácie a výpočtov, ktoré sa môžu udiať iba ak máme k dispozícii nové merania od zariadenia.

\section{Spracovanie sériovej komunikácie}
\begin{lstlisting}[caption=Spracovanie a zaznamenie dát zo sériovej komunikácie., label={code:record}, language=Matlab]
% ----------------------------------
% Process the read data from the serial communication
% ----------------------------------
waitfor(serPort, "UserData");

% Get current time
time_curr = datetime('now');

% Calculate time elapsed since last iteration
time_delta = milliseconds(time_curr - time_tick);

% Read and parse the received data
serLineList = str2num(serPort.UserData); %#ok<ST2NM>

time_tick = time_curr;

% Calculate total time elapsed
time_elapsed = seconds(time_curr - time_start);

% Extract values from the received data
plant_time = serLineList(1) - plant_time_init;
plant_input = serLineList(2);
plant_output = serLineList(3);
plant_potentiometer = R_WANTED + serLineList(4)/100*20;
plant_dt = serLineList(5);

dx = plant_output - timer_yhat(end);
cyhat = timer_yhat(end) + alpha*(dx);
timer_yhat = [timer_yhat, cyhat];
timer_dyhat = [timer_dyhat, timer_dyhat(end) + beta*(dx/time_delta)];

% Record the received data
tmp_printlist = [time_elapsed, plant_time, plant_potentiometer, plant_output, plant_input, plant_dt, time_delta];
doUpdate(tmp_printlist);

% ----------------------------------
% ----------------------------------
\end{lstlisting}

Ako sme spomínali, že samotná slučka čaká na nové dáta z merania, tak to možno vidieť vo výpise \ref{code:record}, kde na začiatku po komentároch voláme funkciu \texttt{waitfor}, ktorá čaká pokiaľ object \texttt{serPort} nezaznamená zmenu v premennej \texttt{UserData}, ktorá sa spätne zapisuje v obslužnej funkcii čítania sériovej zbernice (viď \ref{code:read}). Následne, keď máme k dispozícii nové meranie, tak ako prvé dáta spravujeme a zapíšeme do premenných, ktoré držia hodnotu týchto signálov. Po spracovaní tieto dáta používame na vyhodnotenie $\alpha-\beta$ filtru. Nakoniec všetky signály, ktoré chceme zaznamenať vložíme do premennej \texttt{tmp\_printlist}, čo následne vložíme do funkcie \texttt{doUpdate} ako parameter, aby toto meranie zaznamenala do súboru a vykreslila do grafu. V tejto časti sme definovali aj vlastný program - výpočet filtra, to však z dôvodu toho, že dáta potrebujeme zaznamenať do súboru a grafu, v prípade, že by sme tieto hodnoty nechceli zaznamenať, tak vieme výpočty posunúť do časti vlastného kódu (viď \ref{code:program}), a následne ich použiť pri riadení.

\section{Vlastný program}
\begin{lstlisting}[caption=Blok pre vlastný program, napr. PID regulator., label={code:program}, language=Matlab]
% ----------------------------------
% Insert your code here (for example, we have PID controller code implemented in this space)
% ----------------------------------

e = plant_potentiometer - plant_output;

e_der = (e - e_old) / (time_delta/1000);

e_int = e_int_old + (e * (time_delta/1000));

e_old = e;
e_int_old = e_int;


u = P * e  +  I * e_int + D * e_der;
\end{lstlisting}

Táto časť kódu \ref{code:program} slúži na implementáciu užívateľom navrhnutý riadiaci systém a vlastný kód. Je však potrebné si dávať pozor na to, že nakoľko toto meranie beží v reálny čas, tak dĺžka trvania riadiaceho algoritmu ovplyvní samotné meranie, nakoľko Matlab nebude čakať na dáta zo zariadenia, nakoľko stále trvá výpočet akčného zásahu. Preto je potrebné prispôsobiť periódu vzorkovania dĺžke trvania riadiaceho algoritmu, inak môže dojsť k desynchronizácii dát, či ich stráte. A však výpočet je oveľa rýcheljší ako samotné vykreslovanie, preto v prípade, že potrebujeme čo najpresnejšie meranie pri nízkej perióde vzorkovania, tak pomôže vypnút vykreslovanie v reálnom čase. V konzoli aj v prípade, že si po meraní vykreslíme graf pre premennú \texttt{dt}, tak budeme vidieť stabilnejší priebeh, ako v prípade, že beží aj vykreslovanie.

\section{Saturácia akčného zásahu}
\begin{lstlisting}[caption=Obmedzenie akčného zásahu na maximálne a minimálne hodnoty., label={code:saturation}, language=Matlab]
% ----------------------------------
% Saturate the control output to the MAX and MIN values
% ----------------------------------

u_send = u;

if u_send > U_MAX
    u_send = U_MAX;
elseif u_send < U_MIN
    u_send = U_MIN;
end

% ----------------------------------
% ----------------------------------
\end{lstlisting}

Vnášame do našeho systému nelinearitu (viď \ref{code:saturation}), ktorú nazývame saturácia alebo obmedzenie akčnej veličiný a to len aby sme zabezpečili, že do zariadenia neposielame nevhodné hodnoty vstupných veličín. Týmto zabezpečujeme bezpečný beh zariadenia a najmä bezpečnosť okoliu, minimalizujeme možnosť poškodenia.

\section{Posielanie sériovej komunikácie}
\begin{lstlisting}[caption=Funkcia na posielanie žiadanej akčnej veličiny po sériovej linke., label={code:send}, language=Matlab]
% ----------------------------------
% Send control input to the serial port
% ----------------------------------

write(serPort, u_send, "single");

% ----------------------------------
% ----------------------------------
\end{lstlisting}

Jeden z najdôležitejších riadkov programu, ktorý zabezpečuje odosielanie akčnej veličiny do Arudina na spracovanie a následné prevedenie žiadaného akčného zásahu na zariadenie, môžeme vidieť na výpise \ref{code:send}.

\section{Konečná podmienka merania}
\begin{lstlisting}[caption=Podmienka na bezpečné ukončenie merania., label={code:safety}, language=Matlab]
% ----------------------------------
% Check if the simulation should stop (safety precaution)
% ----------------------------------

if time_elapsed >= T_stop || plant_output >= Y_SAFETY
    configureCallback(serPort, "off"); % Remove the callback from the serial port, before exiting the loop
    break;
end

% ----------------------------------
% ----------------------------------
\end{lstlisting}

Kód v bloku \ref{code:safety} definuje ukočovaciu podmienku merania, v tomto prípade máme dve podmienky ukočenia.

\begin{enumerate}
    \item Keď doba trvania merania prekročí nami nastavenú maximálnu dobu trvania teda \texttt{T\_stop}.
    \item Alebo v prípade, že výstupná veličina zo zariadenia presiahne hodnotu \texttt{Y\_SAFETY}, nakoľko sýstém nevieme ovládať v momente, keď výstupná veličina prevrší hodnotu $180^\circ$.
\end{enumerate}

\section{Ukončenie časovačov}
\begin{lstlisting}[caption=Ukončenie a odstránenie všetkých aktívnych Matlab časovačov., label={code:timer_end}, language=Matlab]
% ----------------------------------
% Close and delete all the existing timers
% ----------------------------------

for tim=timerfindall
    stop(tim);
    delete(tim);
end

% ----------------------------------
% ----------------------------------
\end{lstlisting}

Na výpise \ref{code:timer_end} ukočujeme všetky existujúce časovače v prostredí Matlab, teda aj tie ktoré nespustilo naše meranie. Možno si kód upraviť, tak aby ukočilo iba časovače nami spustené, to však nechávame na čítateľa.

\section{Ukončenie komunikácie}
\begin{lstlisting}[caption=Ukončenie sériovej a súborovej komunikácie, po ukončení merania., label={code:close}, language=Matlab]
% ----------------------------------
% Close the serial connection
% ----------------------------------

% Send a final command and close the serial port
write(serPort, 0.0, 'single');
serPort.flush("input");
clear serPort;
fclose(datafileID);
clear datafileID;

% ----------------------------------
% ----------------------------------
\end{lstlisting}

Zastavujeme zariadenie, ukončujeme sériovú komunikáciu so zariadením a zatvárame súbory, do ktorých sme zapisovali dáta merania. Vymazávame všetky objekty, ktoré sú spojené s týmito akciami, nakoľko sú už po ukončení nepoužitelné. (viď \ref{code:close}) Táto časť programu zabezpečuje bezpečné ukočenie všetkých komunikácií a dovedenie zariadenia do ustáleného stavu.

\section{Uloženie merania}
\begin{lstlisting}[caption=Ukladanie meracích dát do csv a mat súborov., label={code:save}, language=Matlab]
% ----------------------------------
% Save the measurement into a .MAT file for easier access to data when using Matlab
% ----------------------------------

logsout = readtable(FILEPATH, "VariableNamingRule","preserve","Delimiter",",");

save(FILEPATH_MAT);

% ----------------------------------
% ----------------------------------
\end{lstlisting}

Na výpise \ref{code:save} môžeme vidieť spôsob zapisovania dát do súboru \emph{MAT}, kedy ukladáme všetky premenné, ktoré sa nachádzajú v prostredí Matlab. Zároveň ako prvé si možeme všimnúť, akým spôsobom sa majú načítavať merania z \emph{CSV} súboru.

\section{Vykreslenie priebehu merania}
\begin{lstlisting}[caption=Vykreslenie základných veličín procesu merania., label={code:plot_basic}, language=Matlab]
% ----------------------------------
% Quickly plot the measurement - reference, output, and control signal
% ----------------------------------

t = logsout.t;
y = logsout.y;
u = logsout.u;
r = logsout.r;
e = r - y;
dt = logsout.dtp;


figure(111);
hold on;
plot(t, y, '-k', 'LineWidth', 1.5);
plot(t, r, '-r', 'LineWidth', 1.5);
plot(t, u, '-b', 'LineWidth', 1.5);
title('Control Response');
subtitle("P = " + num2str(P) + ", I = " + num2str(I) + ", D = " + num2str(D));
legend('y(t)', 'ref(t)', 'u(t)', "Location", "best");
xlabel('t [s]');
ylabel('y [deg]');
grid on;
hold off;
% ----------------------------------
% ----------------------------------
\end{lstlisting}

Nakoniec merania vo funkcii ešte zobrazíme namerané hodnoty, ktoré sú čítane priamo zo súborov, to slúži na spatnú kontrolu s tým, čo sa vykreslilo na graf reálneho času alebo konečný graf mimo funkcie (viď \ref{code:plot_obs}), ktorý však vykresluje dáta, ktoré sú uložene v pamäti Matlab-u. Preto graf na výpise \ref{code:plot_basic} slúži na spätnú kontrolu.

\section{Vykreslenie priebehu $\alpha-\beta$ filtru}
\begin{lstlisting}[caption=Vykreslenie priebehu a porovnania odhadu stavu pomocou $\alpha-\beta$ filtra., label={code:plot_obs}, language=Matlab]
% ----------------------------------
%% Plot the data
% ----------------------------------

figure(100);
subplot(3, 1, 1);
plot(t, y, t, yhat, t, potentiometer, 'LineWidth', 1.5);
grid minor;
legend('y','yhat','ref');
xlabel('t [s]');
ylabel('$\varphi [^\circ]$', 'Interpreter', 'latex');
title('System response');
subtitle("$\alpha - \beta$ filter", 'Interpreter', 'latex');

subplot(3, 1, 2);
plot(t, dyhat, 'LineWidth', 1.5);
grid minor;
xlabel('t [s]');
ylabel('$\omega [^\circ/s]$', 'Interpreter', 'latex');
title('System velocity response');
subtitle("$\alpha - \beta$ filter", 'Interpreter', 'latex');

subplot(3, 1, 3);
plot(t, (y-yhat), 'LineWidth', 1.5);
grid minor;
xlabel('t [s]');
ylabel('$\varphi [^\circ]$', 'Interpreter', 'latex');
title('Observer error');
subtitle("$\alpha - \beta$ filter", 'Interpreter', 'latex');

% ----------------------------------
% ----------------------------------
\end{lstlisting}

Po ukončení merania vykreslíme dáta \ref{code:plot_obs}, ktoré nám vrátila funkcia v ktorej dané meranie prebiehalo. Toto slúži na ukážku toho, že vieme s dátami pracovať priamo po meraní. V našom prípade vykreslujeme porovnanie nameraných hodnôt a pozorovaných výstupných hodnôt z $\alpha-\beta$ filtra. Tieto grafy sú vytvorené iba z dát, ktoré dostaneme z funkcie \texttt{runArduinoPlot()} a následne spracovanými hodnotami, ktoré sme vytvorili kombináciou meraných veličín.










% -----------------------------------------------------------------------------

\end{document}

% -----------------------------------------------------------------------------